\documentclass[11pt,a4paper]{article}

\usepackage[T1]{fontenc}
\usepackage[utf8]{inputenc}
\usepackage{amsmath,amssymb}
\usepackage[margin=2.7cm]{geometry}

\setlength{\parindent}{0pt}
\setlength{\parskip}{0.5\baselineskip}

\begin{document}

\begin{center}
  {\Large\bfseries Majorization and Uhlmann's Theorem}\\[0.6em]
  Pawel Wocjan\\[0.2em]
  12 July 2002
\end{center}

\vspace{0.8em}

\noindent
The two slides headed \emph{Mathematical tools} from \emph{Complexity of Mutual
Simulation of Hamiltonian Dynamics}, Ludwig-Maximilians-Universität München,
12 July 2002. They are set here on their own because they are the earliest record
I have of my working with majorization, which came back much later in my work on
spectral graph theory. The wording and the notation are those of the 2002 slides;
only the typesetting is new.

\section*{Vector majorization}

Let $x=(x_1,x_2,\ldots,x_d)\in\mathbb{R}^d$.

Define
$x^\downarrow=(x^\downarrow_1,x^\downarrow_2,\ldots,x^\downarrow_d)\in\mathbb{R}^d$
as the vector whose components are the same as those of $x$, but ordered so that
\[
  x^\downarrow_1\ge x^\downarrow_2\ge\ldots\ge x^\downarrow_d\,.
\]

$x$ \emph{majorizes} $y$, written $x\preceq y$, if
\[
  \sum_{j=1}^k x^\downarrow_j \;\le\; \sum_{j=1}^k y^\downarrow_j
\]
for $k=1,\ldots,d-1$ and
\[
  \sum_{j=1}^d x^\downarrow_j \;=\; \sum_{j=1}^d y^\downarrow_j\,.
\]

\section*{Operator majorization}

Let $A$ and $B$ be two Hermitian matrices of size $d$.

We say $A$ \emph{majorizes} $B$, written $A\preceq B$, if
\[
  \lambda(A)\preceq\lambda(B)\,,
\]
where $\lambda(A)$ and $\lambda(B)$ are the spectra (the vectors of eigenvalues)
of $A$ and $B$, respectively.

\subsection*{Uhlmann's theorem}

We have $A\preceq B$ if and only if
\[
  A=\sum_{j=1}^N p_j\, U_j^\dagger B\, U_j\,,
\]
where $\{p_j\}$ is a probability distribution and $\{U_j\}$ is a set of unitaries.

\end{document}
